% informe-entrega-3.tex
%
% Documento maestro del Informe Final de la Tercera Entrega de BIOPET.
% Clase: report. Idioma: espanol (babel). Codificacion: UTF-8.
% Compilacion clasica BibTeX (ver docs/informe/README.md):
%   pdflatex informe-entrega-3.tex
%   bibtex informe-entrega-3
%   pdflatex informe-entrega-3.tex
%   pdflatex informe-entrega-3.tex
% o, si esta disponible, latexmk -pdf -interaction=nonstopmode -halt-on-error informe-entrega-3.tex
%
% Este documento debe compilar aunque ninguna imagen de figuras/ exista
% todavia: cada figura de evidencia esta protegida con \IfFileExists (ver
% el comando \evidencia definido mas abajo), de forma que la ausencia de
% una captura no deja huecos visibles ni texto de marcador en el PDF.

\documentclass[12pt,a4paper]{report}

% --- Idioma y codificacion --------------------------------------------
\usepackage[utf8]{inputenc}
\usepackage[T1]{fontenc}
\usepackage{lmodern}
\usepackage[spanish]{babel}

% --- Geometria de pagina -------------------------------------------------
\usepackage[a4paper,margin=2.5cm]{geometry}

% --- Manejo seguro de guion bajo en modo texto (identificadores de
%     codigo como ROLE_ADMIN, http_status, LOGIN_SUCCESS, etc. se pueden
%     escribir tal cual dentro de \texttt{} sin necesidad de escapar "_")
\usepackage{underscore}

% --- Tablas, figuras, listados --------------------------------------------
\usepackage{graphicx}
\usepackage{float}
\usepackage{booktabs}
\usepackage{longtable}
\usepackage{array}
\usepackage{caption}
\usepackage{enumitem}
\usepackage{textcomp}
\usepackage{amsmath}
\usepackage{amssymb} % \checkmark (tabla CRediT) y simbolos matematicos estandar
\usepackage{pdflscape} % pagina en apaisado solo para la matriz de trazabilidad

% --- Encabezados/pies de pagina -------------------------------------------
\usepackage{fancyhdr}
\pagestyle{fancy}
\fancyhf{}
\fancyhead[L]{BIOPET --- Informe final, Tercera Entrega}
\fancyhead[R]{\leftmark}
\fancyfoot[C]{\thepage}
\renewcommand{\headrulewidth}{0.4pt}

% --- Enlaces internos ------------------------------------------------------
\usepackage[hidelinks]{hyperref}
\hypersetup{
  pdftitle={BIOPET -- Sistema web de gestion veterinaria: Informe final de la Tercera Entrega v0.9.0-rc},
  pdfauthor={Mariscal Cabrera Jaime Josue; Beltran Montiel Fred Adrian; Taipe Mora Zaida Melissa},
  pdfsubject={Informe final de la Tercera Entrega -- Aplicaciones Web},
  bookmarksnumbered=true,
  breaklinks=true,
  colorlinks=false
}

% --- Bibliografia estilo IEEE (numerico), via BibTeX clasico -------------
% "ieeetr" es una hoja de estilo BibTeX estandar (numeros entre corchetes,
% orden de aparicion), disponible en cualquier instalacion completa de
% TeX Live o MiKTeX sin paquetes adicionales.

% --- Comando auxiliar para evidencias condicionales -----------------------
% Uso: \evidencia{ruta/relativa.png}{Caption academico}{fig:label-unico}
% Si el archivo no existe todavia, no se imprime nada (ni hueco, ni texto
% de marcador): \IfFileExists es un comando nativo del nucleo de LaTeX.
\newcommand{\evidencia}[3]{%
  \IfFileExists{#1}{%
    \begin{figure}[H]
      \centering
      \includegraphics[width=0.82\textwidth]{#1}
      \caption{#2}
      \label{#3}
    \end{figure}
  }{}%
}

% --- Metadatos institucionales (usados en la portada) ---------------------
\newcommand{\universidad}{Universidad T\'ecnica Estatal de Quevedo}
\newcommand{\facultad}{Facultad de Ciencias de la Computaci\'on}
\newcommand{\carrera}{Carrera de Software}
\newcommand{\asignatura}{Aplicaciones Web}
\newcommand{\docente}{Guerrero Ulloa Gleiston Ciceron}
\newcommand{\periodoacademico}{2025 -- 2026}
\newcommand{\cursoparalelo}{A}
\newcommand{\equipo}{BMT}
\newcommand{\tituloinforme}{BIOPET --- Sistema web de gesti\'on veterinaria:\\ Informe final de la Tercera Entrega v0.9.0-rc}
\newcommand{\ciudadfecha}{Quevedo, 31 de julio de 2026}

\begin{document}

% ===========================================================================
% PORTADA
% ===========================================================================
\begin{titlepage}
  \centering
  \vspace*{0.5cm}

  {\large \universidad \par}
  {\large \facultad \par}
  \vspace{0.3cm}
  {\normalsize \carrera \par}
  {\normalsize Asignatura: \asignatura \par}

  \vspace{2.2cm}
  {\Huge \bfseries BIOPET \par}
  \vspace{0.4cm}
  {\Large Sistema web de gesti\'on veterinaria \par}
  \vspace{0.6cm}
  {\Large\itshape Informe final de la Tercera Entrega v0.9.0-rc \par}

  \vspace{2.0cm}
  {\large Equipo \textbf{\equipo} \par}
  \vspace{0.6cm}

  {\normalsize
    \begin{tabular}{c}
      Mariscal Cabrera Jaime Josu\'e \\
      Beltr\'an Montiel Fred Adri\'an \\
      Taipe Mora Zaida Melissa \\
    \end{tabular}
    \par
  }

  \vspace{2.2cm}
  {\normalsize Docente: \docente \par}
  \vspace{0.3cm}
  {\normalsize Per\'iodo acad\'emico: \periodoacademico \quad --- \quad Curso/paralelo: \cursoparalelo \par}

  \vfill

  {\normalsize \ciudadfecha \par}
  \vspace{0.5cm}

\end{titlepage}

% ===========================================================================
% MATERIA PRELIMINAR
% ===========================================================================
\pagenumbering{roman}

\tableofcontents
\clearpage
\listoffigures
\clearpage
\listoftables
\clearpage

\pagenumbering{arabic}
\setcounter{page}{1}

% ===========================================================================
% CAPITULOS (uno por archivo de docs/informe/secciones/)
% ===========================================================================
\chapter{Resumen ejecutivo}
\label{cap:resumen-ejecutivo}

\section*{Prop\'osito de BIOPET}
BIOPET es un sistema web de gesti\'on veterinaria desarrollado como Proyecto
Fin de Curso de la asignatura \asignatura. Permite el registro de usuarios,
el inicio de sesi\'on, y la administraci\'on (alta, consulta, actualizaci\'on
y baja l\'ogica) de mascotas asociadas a un due\~no, con control de acceso
diferenciado por rol (\texttt{ROLE_ADMIN}, \texttt{ROLE_VETERINARIO},
\texttt{ROLE_AUXILIAR}, \texttt{ROLE_DUENO}). El sistema no gestiona datos
cl\'inicos reales ni historiales m\'edicos de mascotas: su alcance actual
cubre identificaci\'on, autenticaci\'on y un cat\'alogo b\'asico de mascotas
por due\~no.

\section*{Arquitectura general}
El sistema sigue una arquitectura de tres niveles, documentada formalmente
mediante el modelo C4 (cap\'itulo~\ref{cap:arquitectura-c4}): un
\emph{frontend} Angular~17 (aplicaci\'on de p\'agina \'unica, componentes
\emph{standalone}), un \emph{backend} Spring~Boot~3.2 sobre Java~21 que
expone una API REST, una base de datos relacional PostgreSQL~16 gestionada
con Flyway, y Redis~7 usado tanto para la lista negra de tokens JWT
revocados como para la cach\'e de lecturas de mascotas. Todo el sistema se
levanta mediante Docker Compose y un \texttt{Makefile} con un \'unico
comando reproducible (\texttt{make up}), sin pasos manuales adicionales.

\section*{Principales mejoras de la Tercera Entrega}
Respecto a la Entrega~1B, la Tercera Entrega incorpor\'o, con evidencia
verificable en el repositorio:
\begin{itemize}[nosep]
  \item Migraci\'on completa del flujo de autenticaci\'on de tokens en el
    cuerpo de la respuesta a cookies \texttt{HttpOnly}+\texttt{Secure}+
    \texttt{SameSite=Strict}, con revocaci\'on real v\'ia Redis, rate
    limiting de login, auditor\'ia estructurada y cabeceras de seguridad
    (Jaime; cap\'itulo~\ref{cap:resultados-jaime}).
  \item HTTPS/TLS~1.3 nativo en el backend, con certificado autofirmado
    acad\'emico y verificaci\'on en vivo del protocolo y el cifrado
    negociados (Jaime).
  \item Separaci\'on de privilegios en PostgreSQL (cuenta de aplicaci\'on
    \texttt{biopet_app} sin permisos de DDL), pinning de im\'agenes Docker
    por digest \texttt{sha256}, y medici\'on de rendimiento y cach\'e con
    k6 y Redis (Fred; cap\'itulo~\ref{cap:resultados-fred}).
  \item Integraci\'on del frontend con el nuevo esquema de cookies, CRUD de
    mascotas accesible (atributos ARIA, manejo de errores anunciado a
    lectores de pantalla) y consolidaci\'on de la ingenier\'ia de
    requisitos (SRS, historias de usuario, casos de uso) (Zaida;
    cap\'itulo~\ref{cap:resultados-zaida}).
  \item Cobertura de pruebas verificada autom\'aticamente con JaCoCo
    (umbral m\'inimo del 60\,\% sobre l\'ineas, ramas y complejidad,
    aplicado a todo el m\'odulo del backend).
\end{itemize}

\section*{Seguridad}
El backend fue auditado contra seis categor\'ias aplicables de OWASP~Top~10
(A01, A02, A03, A05, A07, A09), con evidencia reproducible en
\texttt{docs/mediciones/sec/}: autorizaci\'on por rol y por propietario,
HTTPS/TLS~1.3 real, protecci\'on estructural contra inyecci\'on SQL,
cabeceras de seguridad HTTP, autenticaci\'on por cookies con rate limiting y
revocaci\'on, y auditor\'ia de eventos sin datos sensibles. Las seis
categor\'ias auditadas obtuvieron resultado \textsc{pass}, respaldado por
pruebas automatizadas reales (no solo inspecci\'on manual). El detalle
completo, con limitaciones expl\'icitas por categor\'ia, se presenta en el
cap\'itulo~\ref{cap:resultados-jaime}.

\section*{Reproducibilidad}
El sistema completo se reconstruye desde una clonaci\'on limpia con
\texttt{make up}, sin intervenci\'on manual. Las im\'agenes de terceros
(PostgreSQL, Redis, Maven, Eclipse Temurin) est\'an fijadas por digest
\texttt{sha256} para evitar derivas silenciosas entre reconstrucciones. La
inicializaci\'on de la base de datos usa un mecanismo doble documentado en
ADR-004: Flyway como \'unica fuente de verdad del esquema de aplicaci\'on, y
scripts de \texttt{docker-entrypoint-initdb.d} para crear roles con
privilegios m\'inimos y sembrar datos deterministas.

\section*{Rendimiento}
Se ejecutaron tres corridas de k6 en fr\'io y tres en caliente contra el
endpoint cacheado \texttt{GET /api/mascotas}, sobre HTTPS/TLS~1.3 real. La
tasa de error fue del 0{,}0\,\% en las seis corridas, con un throughput
sostenido de aproximadamente 91--92 peticiones por segundo y una mediana de
latencia de 5{,}7 a 9{,}4~ms seg\'un la corrida. El detalle estad\'istico
completo (intervalos de confianza al 95\,\% con la distribuci\'on t de
Student, percentiles p50/p90/p95/p99) se presenta en el
cap\'itulo~\ref{cap:resultados-fred}.

\section*{Accesibilidad y usabilidad}
El frontend incorpora atributos ARIA (\texttt{aria-invalid},
\texttt{aria-describedby}, \texttt{role="alert"}, \texttt{aria-live}) en los
formularios de autenticaci\'on y en el CRUD de mascotas, verificables
directamente en el c\'odigo fuente. \textbf{Las mediciones cuantitativas de
usabilidad (System Usability Scale) y de accesibilidad automatizada
(Lighthouse) todav\'ia no se han ejecutado a la fecha de este informe}: no
existe en el repositorio ni \texttt{docs/mediciones/sus/} ni
\texttt{docs/mediciones/lighthouse/}. Este informe documenta la
metodolog\'ia prevista para ambas mediciones (cap\'itulo~\ref{cap:protocolo-experimental})
sin reportar resultados que a\'un no existen.

\section*{Resultados generales}
La suite de pruebas del backend ejecuta 109 pruebas reales, con 0 fallos, 0
errores y 0 omitidas (\texttt{BUILD SUCCESS}), verificado en esta misma
fase directamente contra los reportes de Surefire y el reporte JaCoCo
local. La cobertura medida fue del 95{,}87\,\% en l\'ineas, 76{,}87\,\% en
ramas y 80{,}00\,\% en complejidad ciclom\'atica, todas por encima del
umbral autom\'atico configurado del 60\,\%.

\section*{Limitaciones principales}
El certificado TLS es autofirmado y exclusivamente acad\'emico; el rate
limiting de login es en memoria y por instancia, no distribuido; los logs
de auditor\'ia no est\'an integrados con un SIEM; la revocaci\'on de tokens
depende de la disponibilidad de Redis; y toda la evaluaci\'on se realiz\'o
sobre una \'unica instancia del backend en un entorno local. Las mediciones
de usabilidad y accesibilidad automatizada quedan pendientes de ejecuci\'on.
Ninguna de estas limitaciones se presenta como resuelta en este informe; se
detallan con su alcance exacto en el cap\'itulo~\ref{cap:amenazas-validez}.

\chapter{Estado actual del sistema}
\label{cap:estado-sistema}

Este cap\'itulo documenta el estado real del sistema BIOPET en la versi\'on
\texttt{v0.9.0-rc}, usando \'unicamente informaci\'on verificable en el
repositorio (c\'odigo fuente, configuraci\'on, documentos de evidencia y
ADR). No se describe ninguna funcionalidad planeada que no est\'e ya
implementada.

\section{Versi\'on y alcance}
La versi\'on etiquetada del repositorio es \texttt{v0.9.0-rc} (candidata a
release de la Tercera Entrega, no una versi\'on de producci\'on). El alcance
funcional cubre autenticaci\'on de usuarios, un CRUD de mascotas con
autorizaci\'on por rol y por propietario, y un endpoint de resumen
estad\'istico por especie.

\section{Frontend}
Aplicaci\'on Angular~17.3 (componentes \emph{standalone}, sin
\texttt{NgModule}), con TypeScript~5.4. Los m\'odulos funcionales
verificados en \texttt{frontend/src/app/} son:
\begin{itemize}[nosep]
  \item \texttt{core/auth.service.ts}, \texttt{core/auth.guard.ts}: sesi\'on
    basada en cookies (sin JWT en \texttt{localStorage}) y protecci\'on de
    rutas.
  \item \texttt{core/http-error.interceptor.ts},
    \texttt{core/problem-detail.service.ts}: interpretaci\'on uniforme de
    errores \texttt{ProblemDetail} del backend.
  \item \texttt{features/login.component.ts}: formulario reactivo con
    validaci\'on y atributos de accesibilidad (\texttt{aria-invalid},
    \texttt{aria-describedby}, \texttt{role="alert"},
    \texttt{aria-live="assertive"}).
  \item \texttt{features/mascotas.component.ts},
    \texttt{features/mascota-api.service.ts}: listado, alta, edici\'on y
    baja de mascotas contra la API real.
\end{itemize}
El detalle funcional y de accesibilidad del frontend se ampl\'ia en el
cap\'itulo~\ref{cap:resultados-zaida}.

\section{Backend}
API REST en Java~21 con Spring~Boot~3.2.12~\cite{spring-boot-docs} (\texttt{Backend/pom.xml}),
organizada en los paquetes \texttt{controller}, \texttt{service},
\texttt{security}, \texttt{repository}, \texttt{exception} y
\texttt{config}. Los tres controladores REST expuestos son
\texttt{AuthController} (\texttt{/api/auth/*}), \texttt{MascotaController}
(\texttt{/api/mascotas/*}) y \texttt{UsuarioController}
(\texttt{/api/usuarios/me}).

\section{PostgreSQL}
PostgreSQL~16, esquema versionado con Flyway
(\texttt{V1__schema_inicial.sql}: tablas \texttt{usuarios} y
\texttt{mascotas}, con \emph{constraint} de rol v\'alido y disparadores de
\texttt{actualizado_en}). Para la reconstrucci\'on reproducible desde
Docker, el esquema se replica en \texttt{db/schema.sql} y la cuenta de
aplicaci\'on (\texttt{biopet_app}) se crea con privilegios m\'inimos en
\texttt{db/roles.sql} (sin \texttt{SUPERUSER}, \texttt{CREATEDB} ni
\texttt{CREATEROLE}), conforme a ADR-004. La funci\'on nativa
\texttt{fn_resumen_mascotas_por_especie} respalda el endpoint de resumen
por especie (cat\'alogo completo en
\texttt{docs/basedatos/CATALOGO-SP.md}).

\section{Redis}
Redis~7, con dos usos reales confirmados en el c\'odigo: lista negra de
JWT revocados (\texttt{TokenBlacklistService}, con TTL igual al tiempo de
vida restante del token) y cach\'e declarativa de lecturas de mascotas
(\texttt{@Cacheable}/\texttt{@CacheEvict} en \texttt{MascotaService}, v\'ia
la abstracci\'on de cach\'e de Spring).

\section{Docker Compose}
Orquestaci\'on de cuatro servicios (\texttt{postgres}, \texttt{redis},
\texttt{backend}, \texttt{frontend}) mediante \texttt{docker-compose.yml},
con un \emph{overlay} opcional \texttt{docker-compose.tls.yml} para activar
el perfil HTTPS. Las im\'agenes de terceros est\'an fijadas por digest
\texttt{sha256}; el arranque completo se realiza con \texttt{make up}
(ADR-005).

\section{HTTPS/TLS}
Perfil Spring \texttt{tls} (\texttt{application-tls.yml}), que activa un
conector HTTPS adicional en el puerto~8443 (TLS~1.3 \'unicamente) junto al
conector HTTP interno en el puerto~8080
(\texttt{TomcatDualConnectorConfig}). El certificado es PKCS12,
autofirmado, generado localmente y nunca versionado.

\section{Autenticaci\'on}
JWT firmado con HMAC-SHA256 (\texttt{jjwt}~0.12.6), transportado en cookies
\texttt{HttpOnly}+\texttt{Secure}+\texttt{SameSite=Strict}
(\texttt{access_token}, \texttt{refresh_token}), con revocaci\'on real v\'ia
Redis, rate limiting de login en memoria y auditor\'ia estructurada. El
detalle completo se documenta en ADR-006 y en el
cap\'itulo~\ref{cap:resultados-jaime}.

\section{Mascotas}
CRUD completo (\texttt{GET}/\texttt{POST}/\texttt{PUT}/\texttt{DELETE}
sobre \texttt{/api/mascotas} y \texttt{/api/mascotas/\{id\}}) con
autorizaci\'on por rol (\texttt{ADMIN}/\texttt{VETERINARIO}/\texttt{AUXILIAR}
para escritura; los cuatro roles para lectura) y por propiedad (un
\texttt{ROLE_DUENO} solo accede a sus propias mascotas). La eliminaci\'on es
l\'ogica (\texttt{activo=false}), no f\'isica.

\section{Cobertura}
Verificada autom\'aticamente con \texttt{jacoco-maven-plugin}~0.8.12
(ejecuci\'on \texttt{check} ligada a la fase \texttt{verify}), con una
regla \texttt{BUNDLE} y un umbral m\'inimo del 60\,\% sobre l\'ineas, ramas
y complejidad. El estado medido actual es \textsc{pass}, con cifras
detalladas en el cap\'itulo~\ref{cap:resultados-jaime}.

\section{Rendimiento}
Medido con k6 contra el endpoint cacheado de mascotas, en tres corridas
fr\'ias y tres calientes, sobre HTTPS/TLS~1.3 real. Tasa de error del
0{,}0\,\% en todas las corridas. Detalle en el
cap\'itulo~\ref{cap:resultados-fred}.

\section{Documentaci\'on}
El repositorio incluye seis ADR (\texttt{docs/adr/}), diagramas C4 de
contenedores y de componentes del backend (\texttt{docs/diagrams/}), un
diagrama entidad-relaci\'on, un diagrama de clases, documentos de evidencia
OWASP y JaCoCo (\texttt{docs/mediciones/sec/}), un reporte de rendimiento
(\texttt{docs/mediciones/perf/}), un diccionario de datos de mediciones
(\texttt{docs/mediciones/DATA-DICTIONARY.md}), una colecci\'on Postman
reproducible (\texttt{docs/postman/}), y documentos de ingenier\'ia de
requisitos (SRS, historias de usuario, casos de uso, matriz de
trazabilidad en \texttt{docs/trazabilidad/matriz.csv}).

\section{Estado de las pruebas}
\begin{table}[H]
  \centering
  \caption{Resultado real de la suite de pruebas del backend (\texttt{mvn clean verify}, reejecutado el 2026-08-01).}
  \label{tab:estado-pruebas}
  \begin{tabular}{@{}ll@{}}
    \toprule
    M\'etrica & Valor \\
    \midrule
    Pruebas ejecutadas & 109 \\
    Fallos & 0 \\
    Errores & 0 \\
    Omitidas & 0 \\
    Cobertura LINE & 95{,}87\,\% \\
    Cobertura BRANCH & 76{,}87\,\% \\
    Cobertura COMPLEXITY & 80{,}00\,\% \\
    Umbral autom\'atico (\texttt{jacoco:check}) & \(\geq\) 60\,\% \\
    Resultado del build & \texttt{BUILD SUCCESS} \\
    \bottomrule
  \end{tabular}
\end{table}

Fuente: \texttt{docs/mediciones/sec/jacoco-summary.md} y verificaci\'on
directa contra \texttt{Backend/target/site/jacoco/jacoco.xml} y
\texttt{Backend/target/surefire-reports/*.txt} generados localmente (no
versionados).

\chapter{Arquitectura C4}
\label{cap:arquitectura-c4}

Este cap\'itulo integra la documentaci\'on arquitect\'onica ya versionada en
\texttt{docs/diagrams/} (Nivel 2 y Nivel 3 del modelo C4~\cite{c4-model}) y las decisiones
de arquitectura relevantes (\texttt{docs/adr/}). Las fuentes de diagrama
son Graphviz (\texttt{.dot}) y PlantUML (\texttt{.puml}) planos, sin macros
de C4-PlantUML, siguiendo la convenci\'on ya establecida en el repositorio.

\section{C4 Nivel 2 --- Contenedores}
\label{sec:c4-nivel-2}

El diagrama de contenedores (\texttt{docs/diagrams/c4-contenedores/})
describe los cuatro contenedores del sistema: el \emph{frontend} Angular, el
\emph{backend} Spring Boot, PostgreSQL y Redis, junto con los protocolos
reales entre ellos (HTTPS/JSON con cookies entre frontend y backend; JDBC
entre backend y PostgreSQL; el cliente Redis de Spring Data entre backend y
Redis). Su fuente ya cuenta con una imagen renderizada en el repositorio
(\texttt{docs/diagrams/c4-contenedores/c4-contenedores.png}).

% EVIDENCIA-COMPARTIDA-C4-L2
% Ruta esperada: figuras/compartidas/c4-contenedores.png
\evidencia{figuras/compartidas/c4-contenedores.png}%
{Diagrama C4 Nivel 2 (contenedores) de BIOPET: frontend, backend, PostgreSQL y Redis, con los protocolos reales entre ellos.}%
{fig:c4-nivel-2}

\section{C4 Nivel 3 --- Componentes del backend}
\label{sec:c4-nivel-3}

El diagrama de componentes del backend
(\texttt{docs/diagrams/c4-componentes-backend/}) fue elaborado mediante
inspecci\'on directa del c\'odigo fuente (constructores, llamadas
directas, configuraci\'on de Spring), no por convenci\'on de nombres.
Documenta 20~componentes reales agrupados en seis \'areas y 4~elementos
externos:

\begin{itemize}[nosep]
  \item \textbf{API REST / Controladores:} \texttt{AuthController},
    \texttt{MascotaController}, \texttt{UsuarioController}.
  \item \textbf{Servicios de aplicaci\'on:} \texttt{AuthService},
    \texttt{MascotaService}.
  \item \textbf{Seguridad y autenticaci\'on:} \texttt{SecurityConfig},
    \texttt{JwtAuthenticationFilter}, \texttt{JwtService},
    \texttt{JwtCookieService}, \texttt{UserDetailsServiceImpl},
    \texttt{LoginRateLimiterService}, \texttt{TokenBlacklistService},
    \texttt{AuthenticationAuditService},
    \texttt{ProblemAuthenticationEntryPoint},
    \texttt{ProblemAccessDeniedHandler}.
  \item \textbf{Manejo de errores:} \texttt{GlobalExceptionHandler},
    \texttt{ProblemDetailFactory}.
  \item \textbf{Persistencia:} \texttt{UsuarioRepository},
    \texttt{MascotaRepository}.
  \item \textbf{Infraestructura del servidor:} \texttt{TomcatDualConnectorConfig}.
  \item \textbf{Elementos externos:} Frontend BIOPET, PostgreSQL, Redis,
    logs locales del proceso.
\end{itemize}

Este diagrama no es un diagrama de clases: no representa atributos,
m\'etodos ni constructores; los DTO y entidades se mencionan \'unicamente
en prosa. Su documento completo, con las siete relaciones principales
(login, autenticaci\'on de una solicitud protegida, consulta/modificaci\'on
de mascotas, logout y revocaci\'on, rate limiting, manejo de errores v\'ia
\texttt{ProblemDetail} y registro de auditor\'ia), est\'a en
\texttt{docs/diagrams/c4-componentes-backend/C4-L3-backend.md}.

\textbf{Nota sobre el renderizado:} la fuente \texttt{c4-componentes-backend.dot}
ya cuenta con una imagen \texttt{.png} generada con Graphviz 15.1.0
(\texttt{dot -Tpng c4-componentes-backend.dot -o c4-componentes-backend.png},
sintaxis validada previamente con \texttt{dot -Tcanon}), disponible en
\texttt{docs/diagrams/c4-componentes-backend/c4-componentes-backend.png}
(5594\texttimes{}2190~px). Falta copiar esa imagen a
\texttt{figuras/jaime/c4-componentes-backend.png} para que aparezca
autom\'aticamente en este informe.

% EVIDENCIA-JAIME-C4-L3
% Ruta esperada: figuras/jaime/c4-componentes-backend.png
\evidencia{figuras/jaime/c4-componentes-backend.png}%
{Diagrama C4 Nivel 3 (componentes) del backend de BIOPET: controladores, servicios, seguridad, manejo de errores, persistencia e infraestructura del servidor.}%
{fig:c4-nivel-3}

\section{Decisiones de arquitectura relevantes (ADR)}
\label{sec:adr-relevantes}

\begin{table}[H]
  \centering
  \caption{ADR vigentes del repositorio y su relaci\'on con la arquitectura descrita.}
  \label{tab:adr-relevantes}
  \begin{tabular}{@{}p{2cm}p{9.5cm}p{2.5cm}@{}}
    \toprule
    ADR & Decisi\'on & Estado \\
    \midrule
    ADR-002 & Migraci\'on de la pila de backend de ASP.NET Core~8 a Java~21 / Spring Boot~3.2. & Aceptado \\
    ADR-003 & Revocaci\'on de JWT mediante lista negra en Redis, indexada por \texttt{jti}. & Aceptado \\
    ADR-004 & Estrategia de base de datos reproducible: Flyway como fuente de verdad del esquema, m\'as \texttt{db/schema.sql}/\texttt{db/roles.sql}/\texttt{db/seed.sql} para inicializaci\'on de Docker con privilegios m\'inimos. & Aceptado \\
    ADR-005 & Reproducibilidad del despliegue: \texttt{Makefile} y pinning de im\'agenes por digest \texttt{sha256}. & Aceptado \\
    ADR-006 & Decisiones de autenticaci\'on y seguridad: JWT por cookies, claims, revocaci\'on, rate limiting, cabeceras, TLS. & Aceptado \\
    \bottomrule
  \end{tabular}
\end{table}

ADR-001 (selecci\'on inicial de ASP.NET Core~8 en la Entrega~1A) qued\'o
formalmente superado por ADR-002 y se conserva \'unicamente como registro
hist\'orico fuera de este repositorio de la Tercera Entrega.

\chapter{Matriz de trazabilidad}
\label{cap:trazabilidad}

La matriz cruda de requisitos vive en
\texttt{docs/trazabilidad/matriz.csv} (id de requisito, tipo, prioridad
MoSCoW, historia de usuario, caso de uso, m\'odulo, endpoint, prueba
autom\'atizada, tipo de acceso, evidencia emp\'irica y estado). Esta secci\'on
selecciona y presenta las filas m\'as representativas de las tres \'areas
de trabajo del equipo, con el estado corregido cuando la evidencia
disponible en el repositorio en el momento de este informe difiere de la
del CSV base.

\textbf{Nota de correcci\'on.} El CSV base marca como \emph{pendiente}
varios requisitos no funcionales de seguridad (TLS, rate limiting,
auditor\'ia) y de rendimiento (cach\'e Redis, k6) porque fue generado antes
de que esas fases de evidencia se fusionaran a \texttt{main}. En esta
matriz esos estados se corrigen a \emph{verificado} \'unicamente cuando
existe un documento de evidencia real y fechado que lo respalda
(\texttt{docs/mediciones/sec/A0X-*.md}, \texttt{docs/mediciones/perf/REPORT.md}).
Ning\'un estado se marca como superado sin esa referencia expl\'icita, y los
requisitos gen\'uinamente pendientes (usabilidad SUS, Lighthouse) se
mantienen como tales.

\begin{landscape}
\begin{longtable}{@{}
  >{\raggedright\arraybackslash}p{1.6cm}
  >{\raggedright\arraybackslash}p{2.0cm}
  >{\raggedright\arraybackslash}p{3.6cm}
  >{\raggedright\arraybackslash}p{3.6cm}
  >{\raggedright\arraybackslash}p{3.6cm}
  >{\raggedright\arraybackslash}p{3.4cm}
  >{\raggedright\arraybackslash}p{3.6cm}
  >{\raggedright\arraybackslash}p{1.8cm}
@{}}
\caption{Matriz de trazabilidad (selecci\'on representativa de \texttt{docs/trazabilidad/matriz.csv} y evidencia posterior).}
\label{tab:trazabilidad} \\
\toprule
Requisito & Responsable & Implementaci\'on & Archivo / componente & Prueba & Medici\'on & Evidencia & Estado \\
\midrule
\endfirsthead
\multicolumn{8}{c}{\tablename\ \thetable\ -- continuaci\'on} \\
\toprule
Requisito & Responsable & Implementaci\'on & Archivo / componente & Prueba & Medici\'on & Evidencia & Estado \\
\midrule
\endhead
\bottomrule
\endfoot

REQ-F-001 & Jaime & Registro de usuario (siempre \texttt{ROLE_DUENO}) & \texttt{AuthController}/\texttt{AuthService} & \texttt{AuthControllerTest} & 109 pruebas, 0 fallos & Postman: Registro (due\~no) & Verificado \\
REQ-F-002 & Jaime & Conflicto de email duplicado (409) & \texttt{GlobalExceptionHandler} & \texttt{AuthControllerTest} & \texttt{ProblemDetail} 409 & \texttt{docs/mediciones/DATA-DICTIONARY.md} & Verificado \\
REQ-F-003 & Jaime & Login por cookies + JWT & \texttt{AuthController}/\texttt{JwtService} & \texttt{AuthControllerTest}, \texttt{JwtServiceTest} & 10 claims JWT & ADR-006 & Verificado \\
REQ-F-004 & Jaime & Refresh de access token & \texttt{AuthController}/\texttt{AuthService} & \texttt{AuthControllerTest} & Cookie renovada & \texttt{A07-authentication.md} & Verificado \\
REQ-F-005 & Jaime & Logout con revocaci\'on Redis & \texttt{AuthController}/\texttt{TokenBlacklistService} & \texttt{AuthControllerTest} & 204 idempotente & \texttt{A07-authentication.md} & Verificado \\
REQ-F-006 & Jaime & Autorizaci\'on por rol (\texttt{@PreAuthorize}) & \texttt{SecurityConfig}/\texttt{MascotaController} & \texttt{MascotaControllerTest} & 401/403 & \texttt{A01-access-control.md} & Verificado \\
REQ-F-007 & Zaida (frontend) / Jaime (endpoint) & Perfil de usuario autenticado & \texttt{UsuarioController} & Postman & \texttt{GET /api/usuarios/me} & \texttt{docs/postman/} & Verificado \\
REQ-F-008--012 & Equipo (base Entrega~1B); accesibilidad: Zaida & CRUD completo de mascotas & \texttt{MascotaController}/\texttt{MascotaService} & \texttt{MascotaControllerTest} & 201/200/204 & \texttt{docs/postman/} & Verificado \\
REQ-F-021 & Fred & Resumen de mascotas por especie (funci\'on nativa) & \texttt{MascotaRepository}, \texttt{fn_resumen_mascotas_por_especie} & \texttt{ResumenEspeciesIntegrationTest} (Testcontainers) & Funci\'on PL/pgSQL & \texttt{docs/basedatos/CATALOGO-SP.md} & Verificado \\
REQ-NF-001/002 & Fred & Cach\'e Redis en listado de mascotas & \texttt{MascotaService} (\texttt{@Cacheable}) & k6 (6 corridas) & Hit ratio cercano al 100\,\% tras calentamiento & \texttt{docs/mediciones/perf/REPORT.md} & Verificado (corrige ``pendiente'' del CSV base) \\
REQ-NF-003 & Jaime & HTTPS/TLS~1.3 nativo & \texttt{TomcatDualConnectorConfig}, \texttt{application-tls.yml} & Verificaci\'on en vivo (\texttt{curl}, \texttt{openssl s\_client}) & TLSv1.3, \texttt{TLS\_AES\_256\_GCM\_SHA384} & \texttt{A02-cryptography-tls.md} & Verificado (corrige ``pendiente'' del CSV base) \\
REQ-NF-004 & Jaime & Generaci\'on/validaci\'on de JWT & \texttt{JwtService} & \texttt{JwtServiceTest} & 10 claims (7 est\'andar + 3 propios) & ADR-006 & Verificado \\
REQ-NF-005 & Jaime & Hash de contrase\~nas (BCrypt) & \texttt{UserDetailsServiceImpl} & Inspecci\'on de BD & Costo BCrypt~12 & \texttt{A07-authentication.md} & Verificado \\
REQ-NF-006 & Jaime & Formato uniforme de error (\texttt{ProblemDetail}) & \texttt{GlobalExceptionHandler}/\texttt{ProblemDetailFactory} & Colecci\'on Postman (40 requests) & RFC~7807 & \texttt{docs/postman/} & Verificado \\
REQ-NF-009 & Jaime & Auditor\'ia estructurada \texttt{AUTH\_AUDIT} & \texttt{AuthenticationAuditService} & \texttt{AuthenticationAuditServiceTest} (10 pruebas) & 7 eventos, INFO/WARN & \texttt{A09-logging.md} & Verificado (corrige ``pendiente'' del CSV base) \\
REQ-NF-010 & Jaime & Rate limiting de login & \texttt{LoginRateLimiterService} & \texttt{AuthControllerTest} (6 pruebas dedicadas) & 5$\times$401, 6.\textsuperscript{o} 429 & \texttt{A07-authentication.md} & Verificado (corrige ``pendiente'' del CSV base) \\
REQ-NF-011 & Fred & Arranque reproducible en un comando & \texttt{Makefile}, \texttt{docker-compose.yml} & \texttt{docker compose ps} (4 servicios \texttt{healthy}) & Digest \texttt{sha256} fijado & ADR-005 & Implementado \\
REQ-NF-012 & Equipo & Estructura de paquetes del backend & \texttt{Backend/src/main/java/com/biopet/**} & Revisi\'on de c\'odigo & --- & C4 Nivel 3 & Verificado \\
REQ-NF-013 & Fred & Privilegios m\'inimos de la cuenta de aplicaci\'on & \texttt{db/roles.sql}, \texttt{biopet_app} & Verificaci\'on manual con \texttt{psql} & Sin \texttt{SUPERUSER}/\texttt{CREATEDB} & ADR-004 & Verificado \\
REQ-NF-007 & Zaida & Accesibilidad/rendimiento del frontend (Lighthouse) & Frontend Angular & Lighthouse CI (pendiente de ejecuci\'on) & --- & --- & Pendiente \\
REQ-NF-008 & Zaida & Usabilidad (SUS) & Frontend Angular & Instrumento SUS, m\'inimo 10 participantes (pendiente) & --- & \texttt{docs/etica/ETHICS.md} & Pendiente \\
\end{longtable}
\end{landscape}

Cobertura de contribuciones reales por integrante en esta matriz: Jaime
(seguridad, autenticaci\'on, autorizaci\'on, TLS, auditor\'ia, formato de
error), Fred (persistencia reproducible, funci\'on nativa de resumen,
cach\'e y rendimiento), y Zaida (accesibilidad del frontend, endpoint de
perfil integrado en la UI, mediciones de usabilidad/accesibilidad a\'un
pendientes de ejecuci\'on). Ning\'un requisito se marca \textsc{pass} o
\emph{verificado} sin una prueba automatizada o un documento de evidencia
citado en la columna correspondiente.

\input{secciones/05-protocolo-experimental}
\chapter{Resultados --- Jaime (seguridad y backend)}
\label{cap:resultados-jaime}

Este cap\'itulo documenta, con evidencia verificable en el c\'odigo y en
\texttt{docs/mediciones/sec/}, el trabajo de seguridad, autenticaci\'on,
pruebas y cobertura desarrollado por Jaime durante la Tercera Entrega. La
auditor\'ia se organiz\'o sobre las categor\'ias aplicables de OWASP Top~10~\cite{owasp-top10}.

\section{Autenticaci\'on JWT por cookies}

El backend emite JWT firmados con HMAC-SHA256 (\texttt{jjwt}~0.12.6). Cada
token contiene \textbf{10 claims}: los \textbf{7 claims est\'andar}
registrados por RFC~7519~\cite{rfc7519} (\texttt{iss}, \texttt{sub}, \texttt{aud},
\texttt{iat}, \texttt{nbf}, \texttt{exp}, \texttt{jti}) y \textbf{3 claims
personalizados} (\texttt{email}, \texttt{rol}, \texttt{typ}). El tipo de
token (\texttt{access} o \texttt{refresh}) se distingue \'unicamente por el
valor del claim \texttt{typ}: ambos comparten el mismo conjunto de nombres
de claims, generados por el mismo m\'etodo (\texttt{buildToken}), y solo
difieren en su tiempo de vida y en ese valor (ADR-006).

Los tokens se transportan en cookies \texttt{access_token} y
\texttt{refresh_token}, ambas con los atributos
\texttt{HttpOnly}+\texttt{Secure}+\texttt{SameSite=Strict}
(\texttt{JwtCookieService}), con rutas diferenciadas
(\texttt{Path=/} para el access token, \texttt{Path=/api/auth} para el
refresh token). El frontend nunca almacena el JWT en \texttt{localStorage};
el navegador env\'ia las cookies autom\'aticamente
(\texttt{credentials: include}). El header \texttt{Authorization: Bearer}
sigue soport\'andose como mecanismo secundario en
\texttt{JwtAuthenticationFilter}, pero el flujo web principal de BIOPET no
lo utiliza.

% EVIDENCIA-JAIME-08
% Captura: atributos de cookies ocultando completamente sus valores.
% Ruta esperada: figuras/jaime/08-cookie-attributes.png
\evidencia{figuras/jaime/08-cookie-attributes.png}%
{Atributos reales de las cookies de sesi\'on (\texttt{HttpOnly}, \texttt{Secure}, \texttt{SameSite=Strict}) inspeccionados en el navegador, sin exponer el valor del token.}%
{fig:jaime-cookie-attributes}

\subsection{Logout y revocaci\'on en Redis}

El logout (\texttt{POST /api/auth/logout}) es idempotente: responde
\texttt{204} incluso sin cookies previas. Cuando hay un token v\'alido, se
extrae su \texttt{jti} y se marca como revocado en Redis
(\texttt{TokenBlacklistService}) con un TTL igual al tiempo de vida
restante del token. En cada solicitud protegida,
\texttt{JwtAuthenticationFilter} consulta esa lista negra \'unicamente para
tokens de tipo \texttt{access}, evitando consultas innecesarias a Redis
para refresh tokens.

\section{Autorizaci\'on por rol y por propietario}

La autorizaci\'on combina dos niveles: por \textbf{rol}
(\texttt{@PreAuthorize} en \texttt{MascotaController}, con
\texttt{ADMIN}/\texttt{VETERINARIO}/\texttt{AUXILIAR} habilitados para
escritura y los cuatro roles habilitados para lectura), y por
\textbf{propiedad} (\texttt{MascotaService.verificarPropiedad}, que lanza
\texttt{AccessDeniedException} cuando un \texttt{ROLE_DUENO} intenta
acceder a una mascota que no le pertenece). Las respuestas se diferencian
correctamente: \textbf{401} cuando no hay autenticaci\'on v\'alida
(\texttt{ProblemAuthenticationEntryPoint}), \textbf{403} cuando hay
autenticaci\'on pero no autorizaci\'on (\texttt{ProblemAccessDeniedHandler}).
Ambos casos producen un \texttt{ProblemDetail} real, sin exposici\'on de
\emph{stack trace}.

% EVIDENCIA-JAIME-05
% Captura: Postman con respuesta 401 ProblemDetail.
% Ruta esperada: figuras/jaime/05-postman-401.png
\evidencia{figuras/jaime/05-postman-401.png}%
{Respuesta 401 (\texttt{ProblemDetail}, \texttt{type=urn:biopet:error:unauthorized}) capturada desde la colecci\'on Postman de BIOPET.}%
{fig:jaime-postman-401}

% EVIDENCIA-JAIME-06
% Captura: Postman con respuesta 403 ProblemDetail.
% Ruta esperada: figuras/jaime/06-postman-403.png
\evidencia{figuras/jaime/06-postman-403.png}%
{Respuesta 403 (\texttt{ProblemDetail}, \texttt{type=urn:biopet:error:forbidden}) capturada desde la colecci\'on Postman de BIOPET.}%
{fig:jaime-postman-403}

\section{Rate limiting de login}

\texttt{LoginRateLimiterService} mantiene, en memoria
(\texttt{ConcurrentHashMap}) y por IP, un contador de intentos fallidos con
ventana y bloqueo de 15~minutos por defecto (\texttt{maxAttempts=6}): los
primeros 5~fallos consecutivos responden \texttt{401}, el 6.\textsuperscript{o}
responde \texttt{429} con cabecera \texttt{Retry-After} (entero positivo de
segundos). El contador se reinicia tras un login exitoso y es
completamente independiente entre IP distintas. Al ser en memoria y por
instancia, no es distribuido entre m\'ultiples r\'eplicas del backend
(limitaci\'on documentada desde la Fase~6).

% EVIDENCIA-JAIME-07
% Captura: sexto intento de login con 429 y Retry-After.
% Ruta esperada: figuras/jaime/07-rate-limit-429.png
\evidencia{figuras/jaime/07-rate-limit-429.png}%
{Sexto intento fallido de login consecutivo: respuesta 429 con cabecera \texttt{Retry-After}.}%
{fig:jaime-rate-limit-429}

\section{\texttt{ProblemDetail} uniforme}

Todos los errores de la API usan el formato RFC~7807~\cite{rfc7807}
(\texttt{application/problem+json}), con los campos \texttt{type},
\texttt{title}, \texttt{status}, \texttt{detail} e \texttt{instance}
(\texttt{GlobalExceptionHandler}, \texttt{ProblemDetailFactory},
\texttt{ProblemType}). El \texttt{detail} nunca incluye \emph{stack trace}
ni el valor bruto de un par\'ametro inv\'alido: para
\texttt{MethodArgumentTypeMismatchException}, el mensaje se construye
\'unicamente con el nombre del par\'ametro (\texttt{ex.getName()}), nunca
con el valor recibido.

\section{Auditor\'ia A09}

\texttt{AuthenticationAuditService} registra siete eventos estructurados en
una sola l\'inea de log (\texttt{AUTH\_AUDIT
timestamp=<UTC> event=<EVENTO> result=<RESULTADO> ip=<IP> subject=<SUJETO>}):
\texttt{LOGIN_SUCCESS}, \texttt{LOGIN_FAILURE}, \texttt{LOGIN_RATE_LIMITED},
\texttt{REFRESH_SUCCESS}, \texttt{REFRESH_FAILURE}, \texttt{LOGOUT_SUCCESS}
y \texttt{TOKEN_REVOKED}, con nivel \texttt{INFO} para los eventos de
resultado esperado y \texttt{WARN} para los fallidos o bloqueados. El
servicio sanitiza \texttt{ip}/\texttt{subject} (elimina caracteres de
control, trunca a 200 caracteres, sustituye valores ausentes por
\texttt{"unknown"}) y \textbf{nunca} recibe como argumento una contrase\~na,
un JWT completo, una cookie ni el encabezado \texttt{Authorization}: la
firma de sus m\'etodos p\'ublicos solo acepta \texttt{ip} y \texttt{subject}.

% EVIDENCIA-JAIME-09
% Captura: log AUTH_AUDIT sin secretos.
% Ruta esperada: figuras/jaime/09-auth-audit.png
\evidencia{figuras/jaime/09-auth-audit.png}%
{L\'inea real de log \texttt{AUTH\_AUDIT} generada durante una prueba de login, sin contrase\~nas ni tokens.}%
{fig:jaime-auth-audit}

\section{Cabeceras de seguridad y TLS}

\texttt{SecurityConfig} (Spring Security~\cite{spring-security-docs}) fija cinco cabeceras en todas las respuestas
(incluidas las p\'ublicas): \texttt{X-Frame-Options: DENY},
\texttt{X-Content-Type-Options: nosniff}, \texttt{Content-Security-Policy:
default-src 'self'; frame-ancestors 'none'; object-src 'none'},
\texttt{Referrer-Policy: no-referrer}, y
\texttt{Strict-Transport-Security} \'unicamente sobre HTTPS
(\texttt{max-age=31536000 ; includeSubDomains ; preload}; ausente sobre
HTTP, comportamiento correcto y verificado).

% EVIDENCIA-JAIME-04
% Captura: curl HTTPS con HSTS y cabeceras de seguridad.
% Ruta esperada: figuras/jaime/04-security-headers.png
\evidencia{figuras/jaime/04-security-headers.png}%
{Cabeceras de seguridad reales capturadas con \texttt{curl.exe -i} contra \texttt{https://localhost:8443/actuator/health}.}%
{fig:jaime-security-headers}

El perfil \texttt{tls} activa un conector HTTPS en el puerto~8443
(\texttt{TLSv1.3} \'unicamente) junto al conector HTTP interno en el
puerto~8080. Verificado en vivo con \texttt{openssl s\_client}: protocolo
\textbf{TLSv1.3}, cifrado \textbf{\texttt{TLS\_AES\_256\_GCM\_SHA384}}
(suite AEAD), certificado autofirmado (\texttt{verify error:num=18:
self-signed certificate}) con SAN \texttt{DNS:localhost,
IP Address:127.0.0.1} --- exclusivamente acad\'emico, nunca v\'alido para
producci\'on.

% EVIDENCIA-JAIME-03
% Captura: OpenSSL mostrando TLSv1.3 y TLS_AES_256_GCM_SHA384.
% Ruta esperada: figuras/jaime/03-tls-openssl.png
\evidencia{figuras/jaime/03-tls-openssl.png}%
{Salida de \texttt{openssl s\_client -connect localhost:8443 -tls1\_3}: protocolo TLSv1.3 y cifrado \texttt{TLS\_AES\_256\_GCM\_SHA384} negociados.}%
{fig:jaime-tls-openssl}

\section{Protecci\'on estructural contra inyecci\'on (A03)}

Todos los m\'etodos de \texttt{UsuarioRepository} y \texttt{MascotaRepository}
son consultas derivadas de Spring Data, con par\'ametros enlazados. La
\'unica \texttt{@Query} nativa del proyecto
(\texttt{fn\_resumen\_mascotas\_por\_especie}) usa un par\'ametro
\texttt{:duenioId} enlazado y fuertemente tipado como \texttt{Long}. Los
payloads de inyecci\'on probados (\texttt{1 OR 1=1},
\texttt{1; DROP TABLE usuarios; --}, \texttt{\%' UNION SELECT NULL --})
nunca llegan a ejecutarse como SQL: son rechazados en la capa de conversi\'on
de Spring MVC (\texttt{MethodArgumentTypeMismatchException}, 400) antes de
alcanzar el repositorio, verificado en
\texttt{SqlInjectionSecurityTest} (incluye la comprobaci\'on de que
\texttt{usuarioRepository.count()} y \texttt{mascotaRepository.count()}
permanecen id\'enticos antes y despu\'es de cada payload).

\section{Postman}

La colecci\'on \texttt{docs/postman/BIOPET.postman\_collection.json}
(Schema~v2.1.0, 40~requests en 5~carpetas) cubre estado del servicio,
autenticaci\'on (incluida la secuencia completa de rate limiting),
mascotas, resumen por especie, y un cat\'alogo de errores controlados
(422, 400, 404, 401, 403, 429). Depende exclusivamente del \emph{cookie
jar} autom\'atico de Postman: ning\'un script guarda un JWT o el valor de
una cookie en una variable, y ning\'un request agrega
\texttt{Authorization: Bearer} al flujo principal.

\section{Cobertura JaCoCo y resultado de pruebas}

JaCoCo~\cite{jacoco-docs} es la herramienta usada para medir y verificar
autom\'aticamente esta cobertura.

\begin{table}[H]
  \centering
  \caption{Cobertura medida por \texttt{jacoco:check} (\texttt{Backend/target/site/jacoco/jacoco.xml}, verificado en esta fase).}
  \label{tab:jacoco-jaime}
  \begin{tabular}{@{}lrrrr@{}}
    \toprule
    Contador & Cubierto & No cubierto & Total & Porcentaje \\
    \midrule
    LINE & 534 & 23 & 557 & 95{,}87\,\% \\
    BRANCH & 103 & 31 & 134 & 76{,}87\,\% \\
    COMPLEXITY & 172 & 43 & 215 & 80{,}00\,\% \\
    \bottomrule
  \end{tabular}
\end{table}

Umbral autom\'atico configurado: \(\geq\) 60\,\% para los tres contadores
anteriores (regla \texttt{BUNDLE} de \texttt{jacoco-maven-plugin}). Las 109
pruebas de la suite completa terminaron con 0~fallos, 0~errores y
0~omitidas (\texttt{BUILD SUCCESS}).

% EVIDENCIA-JAIME-01
% Captura: mvn clean verify con 109 pruebas y BUILD SUCCESS.
% Ruta esperada: figuras/jaime/01-maven-verify.png
\evidencia{figuras/jaime/01-maven-verify.png}%
{Salida real de \texttt{mvn clean verify}: 109 pruebas, 0 fallos, 0 errores, 0 omitidas, \texttt{BUILD SUCCESS}.}%
{fig:jaime-maven-verify}

% EVIDENCIA-JAIME-02
% Captura: resumen HTML de JaCoCo.
% Ruta esperada: figuras/jaime/02-jacoco-resumen.png
\evidencia{figuras/jaime/02-jacoco-resumen.png}%
{Resumen navegable de JaCoCo (\texttt{Backend/target/site/jacoco/index.html}) con la cobertura por paquete.}%
{fig:jaime-jacoco-resumen}

\section{C4 Nivel 3 y ADR-006}

El diagrama de componentes del backend (cap\'itulo~\ref{cap:arquitectura-c4})
y ADR-006 (\emph{Decisiones de autenticaci\'on y seguridad}) documentan
formalmente las decisiones descritas en este cap\'itulo: cookies frente a
almacenamiento en el cliente, revocaci\'on v\'ia Redis, rate limiting en
memoria, y el formato \texttt{ProblemDetail} para todos los errores.

% EVIDENCIA-JAIME-10
% Captura: C4 Nivel 3 renderizado.
% Ruta esperada: figuras/jaime/10-c4-nivel-3.png
\evidencia{figuras/jaime/10-c4-nivel-3.png}%
{Diagrama C4 Nivel 3 (componentes del backend) renderizado, referenciado tambi\'en en la secci\'on~\ref{sec:c4-nivel-3}.}%
{fig:jaime-c4-nivel-3-resultados}

\section{Limitaciones reales de este bloque}

\begin{itemize}[nosep]
  \item El certificado TLS es autofirmado y exclusivamente acad\'emico/local;
    no representa una cadena de confianza v\'alida para producci\'on.
  \item El rate limiting de login es en memoria y por instancia: no es
    distribuido entre m\'ultiples r\'eplicas del backend.
  \item Los logs \texttt{AUTH_AUDIT} se escriben \'unicamente en el log
    local del proceso; no existe integraci\'on con un SIEM centralizado ni
    alertas autom\'aticas.
  \item La revocaci\'on de tokens depende de la disponibilidad de Redis; si
    Redis no est\'a disponible, debe tratarse como error de infraestructura,
    no como sesi\'on v\'alida (ADR-003).
  \item Toda la evidencia de este cap\'itulo se recolect\'o contra
    \textbf{una sola instancia} del backend en un entorno local; no se
    evalu\'o comportamiento bajo m\'ultiples r\'eplicas ni balanceo de
    carga.
\end{itemize}

\input{secciones/07-resultados-fred}
\input{secciones/08-resultados-zaida}
\chapter{Amenazas a la validez}
\label{cap:amenazas-validez}

Este cap\'itulo documenta, sin minimizarlas, las limitaciones metodol\'ogicas
que afectan la interpretaci\'on de los resultados presentados en los
cap\'itulos~\ref{cap:resultados-jaime}--\ref{cap:resultados-zaida}. Ninguna
medici\'on de este informe se presenta como libre de amenazas a su
validez.

\section{Validez interna}

\begin{itemize}[nosep]
  \item Todas las mediciones de seguridad y rendimiento se ejecutaron en
    \textbf{una sola m\'aquina de desarrollo} y contra \textbf{una \'unica
    instancia} del backend; no se control\'o la interferencia de otros
    procesos del sistema operativo durante las corridas.
  \item El certificado TLS es \textbf{autofirmado}, generado localmente
    para esta evaluaci\'on; el protocolo y el cifrado negociados
    (TLSv1.3, \texttt{TLS\_AES\_256\_GCM\_SHA384}) son reales, pero el
    canal de confianza del certificado en s\'i no lo es.
  \item El hit ratio de cach\'e se midi\'o con dos m\'etodos que producen
    cifras distintas ($\approx$49{,}98\,\% global frente a $\approx$100\,\%
    aislado); la elecci\'on del m\'etodo aislado como el m\'as
    representativo es una decisi\'on metodol\'ogica del equipo, documentada
    expl\'icitamente para que un lector externo pueda juzgarla.
  \item El rate limiting de login se evalu\'o \'unicamente en su
    implementaci\'on en memoria, por instancia \'unica; su comportamiento
    bajo m\'ultiples instancias del backend no fue evaluado porque no est\'a
    implementado (no es distribuido).
\end{itemize}

\section{Validez externa}

\begin{itemize}[nosep]
  \item Todas las pruebas se ejecutaron en un \textbf{entorno acad\'emico
    local}, no en un entorno de producci\'on ni de \emph{staging}
    equivalente a producci\'on; los resultados de rendimiento y seguridad
    no deben extrapolarse directamente a un despliegue en la nube con
    m\'ultiples usuarios concurrentes reales.
  \item Los usuarios de prueba (administrador sembrado, cuentas de
    registro) son \textbf{cuentas acad\'emicas de desarrollo}, no usuarios
    reales de una cl\'inica veterinaria; el comportamiento observado del
    sistema frente a datos de producci\'on reales (volumen, variedad,
    calidad de los datos) no fue evaluado.
  \item Lighthouse, cuando se ejecute, depender\'a del entorno de ejecuci\'on
    (hardware, red, versi\'on del navegador) en que se corra: sus
    puntuaciones no son un valor absoluto reproducible en cualquier
    m\'aquina, sino una medici\'on dependiente del entorno.
  \item Un posible \textbf{sesgo de participantes} en la prueba SUS
    (cuando se ejecute) es previsible si los diez participantes m\'inimos
    provienen mayoritariamente del entorno acad\'emico del equipo
    (compa\~neros, familiares con perfil t\'ecnico), en lugar de usuarios
    finales representativos de una cl\'inica veterinaria real.
\end{itemize}

\section{Validez de constructo}

\begin{itemize}[nosep]
  \item Las pruebas de integraci\'on que ejercitan la funci\'on nativa
    \texttt{fn_resumen_mascotas_por_especie}
    (\texttt{ResumenEspeciesIntegrationTest}) usan Testcontainers con
    PostgreSQL real; sin embargo, otras pruebas del backend usan H2 como
    base de datos en memoria bajo el perfil de pruebas, cuyo comportamiento
    \textbf{no es id\'entico} al de PostgreSQL real para funciones nativas
    PL/pgSQL --- por eso el resumen por especie se prob\'o expl\'icitamente
    con Testcontainers en vez de H2.
  \item \texttt{TokenBlacklistService} se prueba siempre con
    \texttt{@MockBean} en las pruebas unitarias/MockMvc; su comportamiento
    real contra una instancia de Redis en condiciones de fallo (ca\'ida de
    red, memoria llena) no se re-verific\'o en esta entrega.
  \item La ausencia de un SIEM centralizado significa que la
    ``auditor\'ia'' medida en este informe es la \textbf{emisi\'on
    correcta} de las l\'ineas de log \texttt{AUTH_AUDIT}, no su
    recolecci\'on, alertamiento o retenci\'on a largo plazo, que son
    aspectos distintos de un programa de auditor\'ia de seguridad.
\end{itemize}

\section{Validez de conclusi\'on}

\begin{itemize}[nosep]
  \item El tama\~no de muestra de cada corrida de k6 ($\sim$3\,200
    peticiones) es adecuado para el c\'alculo de percentiles e intervalos
    de confianza reportados, pero corresponde a una \'unica ventana de
    tiempo corta ($\sim$30--35~s); no se midieron tendencias a lo largo de
    per\'iodos prolongados (horas o d\'ias).
  \item Las 109 pruebas automatizadas del backend y la cobertura JaCoCo
    miden \textbf{cobertura de ejecuci\'on}, no ausencia de defectos: una
    l\'inea cubierta por una prueba no implica que est\'e libre de errores
    l\'ogicos no contemplados por esa prueba.
  \item Algunas evidencias (verificaci\'on manual de privilegios con
    \texttt{psql}, ejecuci\'on manual de \texttt{fn_resumen_mascotas_por_especie})
    se realizaron de forma manual y puntual, no como parte de un pipeline
    de integraci\'on continua que las reejecute autom\'aticamente en cada
    cambio.
\end{itemize}

\input{secciones/10-etica}
\input{secciones/11-credit}
\chapter{Conclusiones}
\label{cap:conclusiones}

\section{Cumplimiento t\'ecnico}
BIOPET \texttt{v0.9.0-rc} implementa y verifica, con evidencia
reproducible, los requisitos funcionales centrales del alcance definido en
el SRS: autenticaci\'on, autorizaci\'on por rol y por propiedad, y un CRUD
completo de mascotas con resumen estad\'istico por especie. La matriz de
trazabilidad (cap\'itulo~\ref{cap:trazabilidad}) muestra la mayor\'ia de los
requisitos funcionales y no funcionales cr\'iticos en estado
\emph{verificado}, con las excepciones expl\'icitas de usabilidad SUS y
accesibilidad autom\'atizada Lighthouse, ambas metodol\'ogicamente
dise\~nadas pero a\'un no ejecutadas.

\section{Seguridad}
Las seis categor\'ias OWASP auditadas (A01, A02, A03, A05, A07, A09)
obtuvieron resultado \textsc{pass}, respaldado por pruebas automatizadas y
verificaci\'on en vivo, no solo por inspecci\'on de c\'odigo. El sistema
implementa autenticaci\'on por cookies seguras, revocaci\'on real de
tokens, rate limiting de login, cabeceras de seguridad HTTP, TLS~1.3 nativo
y auditor\'ia estructurada sin datos sensibles. Estas evidencias no
implican ausencia total de riesgo: las limitaciones de cada control est\'an
documentadas expl\'icitamente en el cap\'itulo~\ref{cap:resultados-jaime} y
en \texttt{docs/mediciones/sec/}.

\section{Reproducibilidad}
El sistema completo se reconstruye desde una clonaci\'on limpia con un
\'unico comando (\texttt{make up}), con im\'agenes de terceros fijadas por
digest para evitar derivas silenciosas, y con separaci\'on de privilegios
verificada en la base de datos. Esto satisface el objetivo del bloque~B de
la gu\'ia de la Tercera Entrega: que un tercero pueda reconstruir el
sistema de forma id\'entica sin intervenci\'on manual.

\section{Rendimiento}
Las seis corridas de k6 realizadas mostraron una tasa de error del
0{,}0\,\% y una latencia mediana de un d\'igito de milisegundos en
condiciones de cach\'e caliente, con un throughput sostenido cercano a
92~peticiones por segundo bajo 50~usuarios virtuales concurrentes. Estos
resultados son alentadores para el alcance evaluado, pero corresponden a
una \'unica m\'aquina de desarrollo y ventanas de tiempo cortas, seg\'un se
detalla en el cap\'itulo~\ref{cap:amenazas-validez}.

\section{Calidad}
La suite de pruebas del backend (109 pruebas, 0 fallos, 0 errores) y la
cobertura verificada autom\'aticamente por JaCoCo (95{,}87\,\% en l\'ineas,
76{,}87\,\% en ramas, 80{,}00\,\% en complejidad, todas por encima del
umbral m\'inimo del 60\,\%) respaldan un nivel de calidad de c\'odigo
consistente con las buenas pr\'acticas exigidas por la gu\'ia. La cobertura
mide ejecuci\'on de pruebas, no ausencia de defectos.

\section{Usabilidad}
El frontend incorpora pr\'acticas de accesibilidad verificables en el
c\'odigo (atributos ARIA, anuncios de error para lectores de pantalla), pero
el equipo no puede todav\'ia afirmar un nivel de usabilidad medido
cuantitativamente: la prueba SUS con participantes externos y la auditor\'ia
Lighthouse quedan como trabajo pendiente, documentado metodol\'ogicamente
en este informe sin resultados inventados.

\section{Limitaciones}
Las limitaciones principales del sistema en su estado actual son: un
certificado TLS autofirmado v\'alido solo para uso acad\'emico; un rate
limiting de login en memoria, no distribuido; aus\'encia de integraci\'on
con un SIEM para los logs de auditor\'ia; dependencia de Redis para la
revocaci\'on de tokens; evaluaci\'on limitada a una sola instancia del
backend; y mediciones de usabilidad/accesibilidad autom\'atica todav\'ia no
ejecutadas.

\section{Estado de \emph{release candidate}}
BIOPET \texttt{v0.9.0-rc} es, como su nombre indica, una \textbf{candidata
a versi\'on}, no una versi\'on de producci\'on. \textbf{El sistema no est\'a
listo para producci\'on}: el certificado TLS es acad\'emico, no existen
mediciones de usabilidad ni de accesibilidad autom\'atizada, la
infraestructura de auditor\'ia carece de un SIEM, y toda la evaluaci\'on de
rendimiento y seguridad se realiz\'o en un entorno local de desarrollo, no
en un entorno equivalente a producci\'on.

\section{Trabajo futuro}
\begin{itemize}[nosep]
  \item Ejecutar la prueba SUS con al menos diez participantes externos al
    equipo y documentar sus resultados agregados en
    \texttt{docs/mediciones/sus/}.
  \item Ejecutar Lighthouse CI sobre el frontend y documentar los
    resultados en \texttt{docs/mediciones/lighthouse/}.
  \item Evaluar un mecanismo de rate limiting distribuido (por ejemplo,
    respaldado en Redis) si el sistema llegara a desplegarse con
    m\'ultiples instancias del backend.
  \item Integrar los logs de auditor\'ia \texttt{AUTH_AUDIT} con un
    sistema de recolecci\'on y alertamiento centralizado (SIEM) antes de
    cualquier despliegue m\'as all\'a del entorno acad\'emico.
  \item Sustituir el certificado autofirmado por uno emitido por una
    autoridad certificadora reconocida si el sistema se despliega fuera
    del entorno acad\'emico.
  \item Completar la migraci\'on de las capturas de evidencia
    (\texttt{figuras/jaime/}, \texttt{figuras/fred/},
    \texttt{figuras/zaida/}, \texttt{figuras/compartidas/}) en la fase
    final de entrega, sin modificar el texto de este informe.
\end{itemize}

\chapter{Anexos}
\label{cap:anexos}

\section{Glosario de siglas}

\begin{table}[H]
  \centering
  \small
  \begin{tabular}{@{}ll@{}}
    \toprule
    Sigla & Significado \\
    \midrule
    ADR & Architecture Decision Record \\
    API & Application Programming Interface \\
    ARIA & Accessible Rich Internet Applications \\
    CRUD & Create, Read, Update, Delete \\
    CSP & Content Security Policy \\
    HSTS & HTTP Strict Transport Security \\
    JWT & JSON Web Token \\
    OWASP & Open Worldwide Application Security Project \\
    RBAC & Role-Based Access Control \\
    SIEM & Security Information and Event Management \\
    SUS & System Usability Scale \\
    TLS & Transport Layer Security \\
    UI & User Interface \\
    WCAG & Web Content Accessibility Guidelines \\
    \bottomrule
  \end{tabular}
  \caption{Siglas usadas en este informe.}
  \label{tab:glosario}
\end{table}

\section{Endpoints principales de la API}

\begin{table}[H]
  \centering
  \small
  \caption{Endpoints reales expuestos por el backend (verificados contra \texttt{AuthController}, \texttt{MascotaController}, \texttt{UsuarioController}).}
  \label{tab:endpoints-anexo}
  \begin{tabular}{@{}llp{7.5cm}@{}}
    \toprule
    M\'etodo & Ruta & Autenticaci\'on / rol \\
    \midrule
    POST & /api/auth/registro & P\'ublico (crea siempre \texttt{ROLE_DUENO}) \\
    POST & /api/auth/login & P\'ublico \\
    POST & /api/auth/refresh & Cookie \texttt{refresh\_token} \\
    POST & /api/auth/logout & Cookie (idempotente sin sesi\'on) \\
    GET & /api/usuarios/me & Cualquier rol autenticado \\
    GET & /api/mascotas & Cualquier rol autenticado (alcance seg\'un rol) \\
    GET & /api/mascotas/\{id\} & Cualquier rol autenticado (propiedad para \texttt{DUENO}) \\
    POST & /api/mascotas & \texttt{ADMIN}/\texttt{VETERINARIO}/\texttt{AUXILIAR} \\
    PUT & /api/mascotas/\{id\} & \texttt{ADMIN}/\texttt{VETERINARIO}/\texttt{AUXILIAR} \\
    DELETE & /api/mascotas/\{id\} & \texttt{ADMIN}/\texttt{VETERINARIO}/\texttt{AUXILIAR} \\
    GET & /api/mascotas/resumen-especies & Cualquier rol autenticado \\
    GET & /actuator/health & P\'ublico \\
    GET & /api/docs & P\'ublico (OpenAPI 3.0, Springdoc) \\
    GET & /api/swagger-ui.html & P\'ublico (documentaci\'on interactiva) \\
    \bottomrule
  \end{tabular}
\end{table}

\section{Documentos de referencia en el repositorio}

\begin{itemize}[nosep]
  \item ADR: \texttt{docs/adr/ADR-002} a \texttt{ADR-007}.
  \item Diagramas C4 y complementarios: \texttt{docs/diagrams/}
    (\texttt{c4-contenedores}, \texttt{c4-componentes-backend},
    \texttt{der-biopet}, \texttt{diagrama-clases}, \texttt{secuencia-jwt}).
  \item DER exportado desde pgAdmin 4 (ERD Tool), evidencia real que cierra
    la observaci\'on OBS-05:
    \texttt{docs/observaciones/evidencias/DER-BIOPET-pgAdmin-ERD-Tool.png}
    (distinto del renderizado Graphviz \texttt{der-biopet.png}, generado
    directamente sobre el esquema real de PostgreSQL).
  \item Evidencia de seguridad: \texttt{docs/mediciones/sec/} (documentos
    \texttt{A01} a \texttt{A09}, \texttt{REPORT.md}, \texttt{jacoco-summary.md}).
  \item Evidencia de rendimiento: \texttt{docs/mediciones/perf/REPORT.md}
    y los archivos \texttt{k6-runN-\{frio,caliente\}.json}.
  \item Diccionario de datos de mediciones:
    \texttt{docs/mediciones/DATA-DICTIONARY.md}.
  \item Colecci\'on y entorno Postman: \texttt{docs/postman/}.
  \item Requisitos: \texttt{docs/requisitos/} (SRS, historias de usuario,
    casos de uso, \texttt{CHANGELOG-REQ.md}) y
    \texttt{docs/trazabilidad/matriz.csv}.
  \item \'Etica y datos: \texttt{docs/etica/ETHICS.md}.
  \item Base de datos: \texttt{docs/basedatos/CATALOGO-SP.md}.
\end{itemize}

\section{Clases de prueba automatizadas citadas en este informe}

\begin{itemize}[nosep]
  \item \texttt{AuthControllerTest}, \texttt{JwtCookieAuthenticationTest},
    \texttt{JwtServiceTest}, \texttt{JwtCookieServiceTest},
    \texttt{LoginRateLimiterServiceTest},
    \texttt{AuthenticationAuditServiceTest} (seguridad y autenticaci\'on).
  \item \texttt{MascotaControllerTest} (autorizaci\'on por rol y por
    propiedad).
  \item \texttt{SqlInjectionSecurityTest}, \texttt{SecurityHeadersTest}
    (OWASP A03 y A05).
  \item \texttt{ResumenEspeciesIntegrationTest} (Testcontainers,
    PostgreSQL real).
  \item \texttt{TomcatDualConnectorConfigTest} (conector TLS adicional).
\end{itemize}

\section{Evidencias pendientes de incorporaci\'on}

Ver la lista completa y numerada de capturas pendientes (responsable,
n\'umero, ruta esperada, descripci\'on) en la secci\'on final del
\texttt{README.md} de este informe (\texttt{docs/informe/README.md}), para
no duplicar esa lista en dos lugares del repositorio.


% ===========================================================================
% BIBLIOGRAFIA (IEEE, BibTeX clasico)
% ===========================================================================
\bibliographystyle{ieeetr}
\bibliography{referencias}

\end{document}
